%==============================================================================
%  Paper 3  --  Elsevier elsarticle format, for Geochimica et Cosmochimica Acta
%
%  Kinetic Arrest, Not Equilibrium:
%  Why Indium Nanoinclusions Survive in Sphalerite
%
%  Compile:  pdflatex paper3_gca.tex   (twice, for cross-references)
%
%  References were checked against the originals by the authors; the
%  markers carried by the working draft have been cleared accordingly.
%==============================================================================
\documentclass[preprint,12pt,authoryear]{elsarticle}

\usepackage{amsmath,amssymb}
\usepackage{booktabs}
\usepackage{graphicx}
\usepackage{lineno}
\usepackage[hidelinks]{hyperref}
\usepackage{url}

\graphicspath{{figures/}{./}}
\journal{Geochimica et Cosmochimica Acta}

\newcommand{\dQ}{\Delta Q}
\newcommand{\JCuIn}{J_{\mathrm{Cu\text{-}In}}}
\newcommand{\Jlike}{J_{\mathrm{like}}}
\newcommand{\Tf}{T_{\mathrm{f}}}

\begin{document}

\begin{frontmatter}

\title{Kinetic Arrest, Not Equilibrium:\\
Why Indium Nanoinclusions Survive in Sphalerite}

\author[hzu]{He Yu}
\ead{yuhe@hzxy.edu.cn}

\author[gut]{Ruqing Chen\corref{cor1}}
\ead{ruqing@hotmail.com}
\cortext[cor1]{Corresponding author.}

\author[uqac]{Huanzhang Lu}
\ead{hzlu@uqac.uquebec.ca}

\address[hzu]{Geological Laboratory, Hezhou University, Hezhou, Guangxi 542899, China}
\address[gut]{GUT Geoservice Inc., Montreal, Quebec, Canada}
\address[uqac]{D\'epartement des sciences appliqu\'ees et Centre d'\'etudes sur
les ressources min\'erales, Universit\'e du Qu\'ebec \`a Chicoutimi, Quebec, Canada}

\begin{abstract}
\noindent
The equilibrium state of indium in sphalerite is a single condensed domain: the
preceding paper of this series established this across every composition from
0.10 to 2.00~at.\%, and we show here that dividing one domain into two costs
between 300 and 1800~$k_BT$, from an interfacial energy obtained in closed form,
$\gamma = \lambda/3 - J_{\rm Cu\text{-}In}/3 - J_{\rm like}/6$, in which the bulk
terms cancel identically. Natural sphalerite nonetheless carries populations of
discrete nanometric and micrometric inclusions. The observed configurations lie
hundreds of $k_BT$ above the state the free energy selects.

We argue that this is kinetic arrest rather than a failure of the
thermodynamics, and that the argument converts a spatial question --- how large
are the inclusions --- into a temporal one. Coarsening is limited by solute
transport, which is thermally activated, so its timescale grows exponentially as
temperature falls; the cooling timescale falls only as $T^2$. The two must
cross, and cross once. Below the crossing the microstructure ceases to evolve
and what survives is the configuration the system occupied when transport
stopped.

The arrest condition is strongly buffered: every quantity entering it, including
the diffusion prefactor this work cannot determine, appears inside a logarithm,
so an order-of-magnitude error shifts the arrest temperature by less than a
quarter. The frozen domain size consequently follows a power law in cooling
rate, $R_{\rm f} \propto |\dot{T}|^{-n}$, whose exponent rather than whose
prefactor carries the physics.

Two calculations support this, with roles that are not interchangeable. The bulk
free energy is measured by thermodynamic integration rather than assumed,
because a regular-solution form has no driving force in this system at all --- the
heterovalent and homovalent contributions cancel identically under random
mixing --- and the gradient energy coefficient is then fixed by $\gamma$. A
phase-field calculation on that footing reproduces the sequence from
homogeneous solution through spinodal decomposition to coarsening and arrest,
and gives $n = 0.259$; it cannot give an absolute length, because the interface
must be numerically widened for any solver to resolve it. Direct lattice
coarsening, with the interface at its physical width of one bond, gives domains
of one to two nanometres and $n = 0.150$.

Set against deposit classes whose cooling rates span eight orders of magnitude,
from quenched seafloor systems to basinal brines, the model reproduces the
observed ordering of inclusion size but implies an exponent near $0.45$, larger
than either simulation. We report the disagreement rather than the agreement.
It may reflect a crossover from interface-limited coarsening at the atomic scale
towards diffusion-limited coarsening as domains grow, in which case no single
exponent describes the whole range. Until it is resolved the inversion supports
ranking deposits by cooling rate, not calibrating one.

Two consequences follow regardless. A smooth LA-ICP-MS depth profile is
evidence of rapid cooling rather than of solid solution, since the indium is
aggregated but below the resolution of the ablation volume. And indium
deportment, and hence its metallurgical recovery, should differ systematically
between deposit classes along an axis predictable from thermal history.
\end{abstract}

\begin{keyword}
sphalerite \sep indium \sep kinetic arrest \sep coarsening \sep phase field \sep
Cahn--Hilliard \sep nanoinclusions \sep cooling rate \sep ore deposits \sep
critical metals
\end{keyword}

\end{frontmatter}

\noindent\small
\textbf{Version 2.0.} This version reports the outcome of a test that version 1
proposed but had not carried out. Section~\ref{sec:crosscheck} noted that four
determinations of the coarsening exponent span a factor of three, and
Section~\ref{sec:discussion} suggested this might reflect a crossover from
interface-limited coarsening at the atomic scale towards diffusion-limited
coarsening as domains grow, adding that the suggestion was testable within a
single simulation. Phase-field runs at four interface widths have since been
made and the hypothesis is \emph{not} supported: the curves do not collapse
when rescaled by interface width, and the exponent rises rather than falls. The
exponent instead tracks how far each run has progressed, at a correlation of
0.99 with the logarithm of the step count, which points to a common asymptote
approached from below. No other conclusion of version 1 is altered; the
discrepancy of Section~\ref{sec:geospeedometer} returns to the unpaired
end-member data as the remaining explanation. Detail in
Section~\ref{sec:discussion}.
\normalsize

\linenumbers


\section{Introduction}
\label{sec:intro}

\subsection{A result and a contradiction}

The preceding paper of this series established, by replica-exchange sampling of
a configurational Hamiltonian whose charge-compensation coefficient is fixed by
the dielectric constant of ZnS rather than fitted, that the equilibrium state of
indium in sphalerite is condensed \citep{yu2026b}. Across every composition
examined, from 0.10 to 2.00~at.\%, the solute population gathered into a
\emph{single} domain: the monomer fraction was zero throughout, against a
random-solution expectation of 98\% at the most dilute point, and the result was
established by bracketing rather than by convergence, so it does not rest on
equilibrating a sampler that cannot be equilibrated at ore-forming temperatures.

Section~\ref{sec:gamma} of the present paper quantifies why the state is
singular rather than merely condensed. The interfacial energy has a closed form,
$\gamma = \lambda/3 - \JCuIn/3 - \Jlike/6$, in which the bulk terms cancel
identically; dividing one domain into two costs between 300 and 1800~$k_BT$ over
the size range examined. Bulk configurational thermodynamics opposes a
population of separate domains by an enormous margin.

Natural sphalerite frequently contains one. High-resolution microscopy resolves
indium-bearing heterogeneity at the scale of tens of nanometres in samples whose
LA-ICP-MS depth profiles are perfectly smooth \citep{cook2009,frenzel2016}, and
roquesite is reported as discrete micrometric grains exsolved from sphalerite
\citep{moura2007}. Neither texture is a single domain. The contradiction is
sharp, and it is quantitative rather than a matter of interpretation: the
observed configurations lie hundreds of $k_BT$ above the state the free energy
selects.

\subsection{The contradiction is only apparent}

Nothing in the preceding result asserts that a sphalerite grain \emph{attains}
its equilibrium configuration. Equilibrium thermodynamics identifies the state a
system would reach given unlimited time; it says nothing about whether the time
was available. An ore body cools, and as it cools the solute transport that
would carry the microstructure towards equilibrium slows exponentially, while
the cooling itself proceeds at a rate set by the geology. Below some
temperature the second outruns the first and the microstructure stops evolving.
What survives to be observed is not the equilibrium state but the state the
system happened to occupy when transport ceased.

This is kinetic arrest, and the proposal of this paper is that it accounts for
the discrepancy in full. The multi-domain textures of natural sphalerite are
neither an alternative equilibrium nor evidence against the thermodynamics of
\citet{yu2026b}; they are that thermodynamics interrupted.

The proposal has a testable consequence, and it is the reason the argument is
worth making rather than merely asserting. If inclusions are frozen coarsening
products, their size records how long coarsening had before it stopped, and
therefore how fast the rock cooled. The spatial question --- how large are the
inclusions --- becomes a temporal one.

\subsection{What is required to make the argument}

The transformation from a spatial question to a temporal one requires three
things, and this paper supplies them in order.

First, the thermodynamics of the lattice model must be carried into a continuum
description in which coarsening can be followed. Section~\ref{sec:bridge} does
this without adjustable parameters: the bulk free energy is measured by
thermodynamic integration rather than assumed, because a regular-solution form
has \emph{no} driving force in this system --- the heterovalent and homovalent
contributions cancel identically under random mixing --- and the gradient energy
coefficient is fixed by the interfacial energy, which Section~\ref{sec:gamma}
obtains in closed form and verifies against explicit construction.

Second, the competition between coarsening and cooling must be made explicit.
Section~\ref{sec:timescales} shows that the coarsening time grows exponentially
as temperature falls while the cooling time falls only as $T^2$, so the two must
cross, once; that the crossing temperature is buffered logarithmically against
every quantity entering it, including the diffusion prefactor this work cannot
determine; and that the frozen domain size consequently follows a power law in
cooling rate.

Third, the prediction must be confronted with observation.
Section~\ref{sec:geospeedometer} sets the power law against deposit classes
whose cooling rates differ by eight orders of magnitude and whose inclusion
textures differ by four. The ordering is reproduced. The magnitude is not, and
we report the disagreement and its possible origins rather than the agreement.

\subsection{Two methods, and what each is for}

Two calculations are used and their roles are not interchangeable, a division
set out fully in Section~\ref{sec:kappacaveat} and stated here because it
governs how every subsequent figure should be read.

A phase-field calculation supplies the pathway and the topology: which
morphologies the system passes through, whether coarsening completes or
arrests, and how the outcome depends on the ratio of cooling rate to mobility.
It cannot supply an absolute length, because the Cahn--Hilliard interface must
be numerically widened for any solver to resolve it, so the domain sizes in
those figures are set by the grid rather than predicted.

A direct lattice coarsening simulation supplies the absolute scale, with the
interface at its physical width of one bond and no inflation anywhere. It
reaches far less time and far fewer domains, but its nanometres are nanometres.

Where both can compute a quantity --- the coarsening exponent --- they are
compared, and Section~\ref{sec:crosscheck} reports that they disagree by 70\%.
That disagreement, and what it does and does not undermine, is treated there
rather than deferred.

%==============================================================================
\section{From lattice to continuum}
\label{sec:bridge}

The Cahn--Hilliard equation \citep{cahn1958} for a conserved order parameter
$c(\mathbf{r}, t)$,
here the total solute fraction $x_{\rm Cu} + x_{\rm In}$ on the cation
sublattice, is
\begin{equation}
\frac{\partial c}{\partial t}
  = \nabla \cdot \left[ M(T)\, \nabla \frac{\delta F}{\delta c} \right],
\qquad
F[c] = \int \left[ f(c) + \tfrac{1}{2}\kappa \left|\nabla c\right|^2 \right] \mathrm{d}V .
\label{eq:ch}
\end{equation}
Three quantities must be supplied: the bulk free energy density $f(c)$, the
gradient energy coefficient $\kappa$, and the mobility $M(T)$. This section
obtains the first two from the lattice model without adjustable parameters, and
sets out what is and is not known about the third.

\subsection{Why $f(c)$ cannot be assumed}
\label{sec:whynotregular}

The usual practice is to postulate a regular-solution or double-well $f(c)$ and
fit its coefficients. That is not available here, and the obstruction is
specific to this Hamiltonian rather than a matter of accuracy.

Under random mixing at $x_{\rm Cu} = x_{\rm In} = c/2$, the configurational part
of the charge-compensation energy averages to
\begin{equation}
\Big\langle \sum_{\langle ij \rangle} \delta q_i \delta q_j \Big\rangle_{\rm rand}
= \frac{Nz}{2}\left[ 2\left(\frac{c}{2}\right)^2(-1)
                   + \left(\frac{c}{2}\right)^2(+1)
                   + \left(\frac{c}{2}\right)^2(+1) \right] = 0,
\label{eq:mfzero}
\end{equation}
because the heterovalent Cu--In contacts, carrying $\delta q_i \delta q_j = -1$,
exactly cancel the homovalent Cu--Cu and In--In contacts, carrying $+1$. A
mean-field free energy built on random mixing therefore has \emph{no driving
force whatever}, at any composition. The entire thermodynamics of this system
resides in the correlations that mean-field theory discards, which is the same
statement, in the language of free energy, as the observation in
\citet{yu2026b} that the equilibrium state is condensed while the random solid
solution is not.

The consequence is that a regular-solution $f(c)$ would not be an approximation
to this system; it would describe a different one. $f(c)$ must be measured.

\subsection{Measuring $f(c)$ by thermodynamic integration}
\label{sec:fofc}

We obtain $f(c)$ by integrating over the coupling strength. Introducing a
scaling parameter $a \in [0,1]$ and writing $H(a) = aH$, the free energy of the
fully coupled system is
\begin{equation}
F(1) = F(0) + \int_0^1 \langle H \rangle_a \, \mathrm{d}a,
\label{eq:ti}
\end{equation}
where $F(0)$ is the ideal-solution free energy, known analytically as
$-k_BT \sum_\alpha x_\alpha \ln x_\alpha$ per site, and $\langle H \rangle_a$ is
the expectation of the full Hamiltonian in the ensemble generated by $H(a)$.
The identity is exact; the only error is the sampling error in
$\langle H \rangle_a$ and the quadrature error, both of which are controlled.

Scaling the \emph{whole} Hamiltonian rather than one term matters. Integrating
along a path that scales only $E_c$, say, would give the free energy difference
along that path, which is a different quantity; the identity of
Eq.~\eqref{eq:ti} requires that $a$ multiply $H$ in its entirety.

Table~\ref{tab:fofc} gives the result at 573~K, computed in an $L = 12$ box
($N = 6912$ cation sites) with nine Gauss--Legendre nodes in $a$ and
block-averaged sampling at each node.

\begin{table}[htbp]
\centering
\caption{Bulk free energy per cation site at 573~K, by thermodynamic
integration. $f_{\rm ideal}$ is the ideal-solution term and the integral column
is $\int_0^1 \langle H\rangle_a\,\mathrm{d}a$. Uncertainties are propagated from
the block-averaged sampling error at each quadrature node.}
\label{tab:fofc}
\small
\begin{tabular}{rrrrr}
\toprule
$c$ & $f_{\rm ideal}$ (eV) & integral (eV) & $f(c)$ (eV) & $\pm$ \\
\midrule
0.005 & $-0.00173$ & $0.00232$ & $+0.00059$ & 0.00002 \\
0.010 & $-0.00311$ & $0.00399$ & $+0.00089$ & 0.00003 \\
0.020 & $-0.00553$ & $0.00653$ & $+0.00101$ & 0.00004 \\
0.040 & $-0.00966$ & $0.01049$ & $+0.00083$ & 0.00024 \\
0.080 & $-0.01650$ & $0.01436$ & $-0.00214$ & 0.00038 \\
0.120 & $-0.02222$ & $0.01965$ & $-0.00258$ & 0.00045 \\
0.160 & $-0.02719$ & $0.02313$ & $-0.00405$ & 0.00041 \\
0.200 & $-0.03155$ & $0.02642$ & $-0.00513$ & 0.00052 \\
\bottomrule
\end{tabular}
\end{table}

\subsection{The instability, and how to diagnose it robustly}
\label{sec:spinodal}

The Cahn--Hilliard instability requires $f''(c) < 0$. Extracting a second
derivative from eight measured points is, however, an ill-conditioned operation,
and we report it with the corresponding caution.

Fitting a polynomial and differentiating twice gives answers that depend on the
degree chosen: at $c = 0.20$ the third-, fourth- and fifth-degree fits return
$f'' = +0.62$, $-2.90$ and $+7.23$~eV respectively. Those are not three
estimates of one number; they are the oscillation of an under-determined fit.
Finite differences on the raw points are better behaved but still change sign
between adjacent intervals in the concentrated range.

The physically meaningful test is not local curvature but the common-tangent
construction, which asks whether a homogeneous state at composition $c$ has
higher free energy than a mixture of two other compositions, and which requires
no differentiation at all. Taking the chord between the end members of the
measured range, the excess $f(c) - f_{\rm chord}(c)$ is
\begin{center}
\small
\begin{tabular}{rrrrrrrr}
\toprule
$c$ & 0.005 & 0.010 & 0.020 & 0.040 & 0.080 & 0.120 & 0.160 \\
\midrule
$f - f_{\rm chord}$ ($10^{-5}$ eV) & $0$ & $+44$ & $+86$ & $+126$ & $-53$ & $+21$ & $-9$ \\
\bottomrule
\end{tabular}
\end{center}
The dilute compositions from 0.010 to 0.040 lie above the chord by three to
five times their sampling uncertainty and are therefore unstable to phase
separation. The concentrated compositions lie on or below it to within
uncertainty.

Two points follow. First, the instability is located in the dilute range,
which is where \citet{yu2026b} found complete condensation and where natural
sphalerite sits: the free-energy measurement and the direct simulation agree,
by independent routes. Second, resolving the phase boundary itself would require
finer composition spacing and smaller error bars than the present sampling
provides; we use $f(c)$ as the driving force in Eq.~\eqref{eq:ch} and do not
claim to have located the binodal.

The distinction between the two kinds of instability matters for what follows.
A composition lying above the common tangent but at positive curvature is
metastable: it separates only by nucleation, over a barrier. One at negative
curvature is unstable outright and separates spontaneously, at every wavelength
longer than a critical value --- spinodal decomposition. The simulations of
Section~\ref{sec:pf} are conducted in the latter regime, so that the pathway is
displayed without the additional and separately uncertain physics of nucleation
rates. Whether natural sphalerite decomposes spinodally or by nucleation depends
on where its composition sits relative to a boundary this work does not resolve;
the arrest physics that the paper concerns is the same either way, because it
turns on transport ceasing rather than on how decomposition began.

\subsection{The gradient coefficient is fixed by the interfacial energy}
\label{sec:kappa}

For a planar interface at equilibrium, the Cahn--Hilliard functional gives the
interfacial energy as
\begin{equation}
\gamma = \int_{c_\alpha}^{c_\beta} \sqrt{2\kappa\, \Delta f(c)}\ \mathrm{d}c,
\qquad
\Delta f(c) = f(c) - f_{\rm tangent}(c),
\label{eq:gammakappa}
\end{equation}
where $f_{\rm tangent}$ is the common tangent between the coexisting
compositions $c_\alpha$ and $c_\beta$. Since $\gamma$ is known in closed form
from the lattice model (Section~\ref{sec:gamma}) and $f(c)$ is now measured,
Eq.~\eqref{eq:gammakappa} determines $\kappa$. It is not a fitting parameter.

\subsection{What this bridge cannot deliver: the absolute length scale}
\label{sec:kappacaveat}

This subsection states a limitation that governs how the phase-field results of
Section~\ref{sec:pf} should be read, and we state it before presenting any of
them.

The characteristic length of the Cahn--Hilliard equation is
$\ell = \sqrt{\kappa / |f''|}$, and it sets the interface width. With $\kappa$
fixed by Eq.~\eqref{eq:gammakappa} from a $\gamma$ that is an atomic-scale
quantity --- the lattice model's interface is one bond wide --- the physical
$\ell$ is of order the lattice parameter. No finite-difference or spectral
solver can represent an interface that thin: numerical stability requires the
interface to span four to six grid points, so any practical calculation must
either use a grid spacing far below the lattice parameter, which is
meaningless, or inflate $\kappa$ so that the interface spreads across the
available grid.

Every phase-field calculation reported here takes the second course, as every
phase-field calculation of an atomically sharp interface must. The consequence
is unavoidable and must not be obscured: \emph{the domain sizes appearing in the
phase-field simulations are set by the grid spacing, not predicted by the
physics}. A reader who asks why the simulated inclusions are nanometric would be
answered, truthfully, that the grid was chosen that way.

We therefore divide the labour between the two methods explicitly.

\begin{itemize}\itemsep2pt
\item \textbf{The phase-field model supplies the pathway and the topology}: the
sequence of morphologies through which the system passes, whether coarsening
completes or arrests, and how that outcome depends on the ratio of cooling rate
to mobility. These are dimensionless statements and are unaffected by the
inflation of $\kappa$, because inflating $\kappa$ rescales lengths and times
together.
\item \textbf{The lattice model supplies the absolute scale}: domain sizes in
atoms and in nanometres, obtained by direct coarsening simulation
(Section~\ref{sec:latticecoarsen}) in which the interface is what it physically
is and no inflation occurs.
\end{itemize}

The two are cross-checked in Section~\ref{sec:crosscheck} on the quantity they
can both compute, namely the coarsening exponent.

\subsection{Mobility}
\label{sec:mobility}

The mobility is the one input this work cannot obtain from the lattice model,
because Monte Carlo sweeps are not physical time. We therefore build the
framework in terms of a dimensionless mobility and state separately what is
required to give it absolute units.

Solute transport in a substitutional sublattice is thermally activated, so
\begin{equation}
M(T) = \frac{D(T)\, \Omega}{k_B T}, \qquad
D(T) = D_0 \exp\!\left(-\frac{E_a}{k_B T}\right),
\label{eq:arrhenius}
\end{equation}
with $\Omega$ the volume per cation site. The activation energy $E_a$ has a
lower bound from the model itself: moving a solute out of a domain requires
breaking its heterovalent bonds, each costing $2\lambda \approx 0.6$~eV, so
$E_a \gtrsim 2\lambda$ for a surface solute and several times that for an
interior one. This is a bound, not a determination: the true $E_a$ also contains
the migration barrier of the underlying diffusion mechanism, whether vacancy- or
interstitial-mediated, which the configurational model does not represent.

Two regimes of conclusion follow, and we label them throughout.

\emph{Conclusions requiring only the Arrhenius form.} The competition between
coarsening and cooling developed in Section~\ref{sec:timescales}, the existence
of an arrest temperature, and the scaling of final domain size with cooling rate
all follow from the \emph{form} of Eq.~\eqref{eq:arrhenius} together with
$E_a > 0$. They are expressed in reduced variables and require no absolute
value of $D_0$.

\emph{Conclusions requiring $D_0$ and $E_a$.} Placing the arrest temperature on
an absolute scale, and converting an observed domain-size distribution into a
cooling rate in kelvin per year, both require measured diffusion data for Cu or
In in sphalerite. Where such values are available they are used and identified
as such; where they are not, results are reported as relative calibrations, in
which a domain size ratio maps onto a cooling rate ratio without either being
absolute. We regard the relative calibration as the safer claim and the one the
present work establishes.

%==============================================================================
\section{Interfacial energy and the splitting barrier}
\label{sec:gamma}

The interfacial energy between a solute domain and the ZnS matrix is required
twice over: to fix $\kappa$ through Eq.~\eqref{eq:gammakappa}, and to set the
thermodynamic penalty for the multi-domain configurations whose survival this
paper explains. It can be obtained in closed form.

\subsection{Closed form}
\label{sec:gammaclosed}

Consider a domain of $n$ solutes, half Cu and half In, and let $A$ be the number
of solute--Zn first-neighbour bonds, which is the interface measured in the
natural units of the lattice. The solutes present $12n$ bond ends, of which $A$
reach the matrix, so the solute--solute bond count is $N_{ss} = (12n - A)/2$. In
an ideal two-Cu-two-In ordering each cation has eight opposite and four like
neighbours, so two thirds of the solute-directed ends are heterovalent:
\begin{equation}
N_{\rm CuIn} = \tfrac{1}{3}(12n - A) = 4n - \tfrac{A}{3},
\qquad
N_{\rm like} = 2n - \tfrac{A}{6}.
\label{eq:bondcounts}
\end{equation}
Substituting into the exact pair form of the charge-compensation energy
\citep{yu2026},
\begin{equation}
E_c = 4\lambda \sum_i \delta q_i^{\,2}
    + 2\lambda \sum_{\langle ij \rangle} \delta q_i \delta q_j
    = 16\lambda n - \lambda A - 4\lambda\left(4n - \frac{A}{3}\right)
    = \frac{\lambda}{3} A .
\label{eq:Ecgamma}
\end{equation}
The bulk terms cancel identically, leaving an energy strictly proportional to
the interface area. Adding the chemical term gives
\begin{equation}
\boxed{\;
\gamma = \frac{\lambda}{3} - \frac{\JCuIn}{3} - \frac{\Jlike}{6}
\;}
\label{eq:gamma}
\end{equation}
per solute--Zn bond, with the corresponding bulk coefficient
\begin{equation}
e_{\rm bulk} = 4\JCuIn + 2\Jlike
             + \tfrac{1}{2}\left(E_s^{\rm Cu} + E_s^{\rm In}\right)
\label{eq:ebulk}
\end{equation}
per solute. The Eshelby term is the \emph{mean} over the two species, not the
value for indium alone, because $e_{\rm bulk}$ is expressed per solute atom and
a stoichiometric domain contains equal numbers of each. Copper contributes
nothing to it: $r_{\mathrm{Cu^+}} = r_{\mathrm{Zn^{2+}}} = 0.60$~\AA{} in
fourfold coordination, so its misfit vanishes identically, and
$E_s^{\rm Cu} = 0$ while $E_s^{\rm In} = 0.033$~eV. At the parameters of
\citet{yu2026}, $\gamma = 0.125$~eV per bond and $e_{\rm bulk} = -0.2835$~eV per
solute.

Two features of Eq.~\eqref{eq:gamma} deserve emphasis. It contains no fitted
quantity: $\lambda$ is fixed by the dielectric constant of ZnS, and the $J$
terms enter explicitly, so the expression updates immediately for any future
first-principles determination without repeating a simulation. And the
cancellation in Eq.~\eqref{eq:Ecgamma} is exact rather than asymptotic, so the
result holds for domains of any size, not only large ones.

\subsection{Verification by explicit construction}
\label{sec:gammaverify}

Domains of seven sizes from 64 to 3072 solutes were built explicitly in an
$L = 24$ box ($N = 55{,}296$ cation sites) and evaluated exactly, then relaxed
by solute--solute exchange alone at 20~K. Restricting the move set in that way
lets the interface reconstruct chemically while holding the domain outline, and
therefore $A$, exactly fixed; it separates the interfacial term from the
dissolution physics.

\begin{table}[htbp]
\centering
\caption{Interfacial and bulk coefficients from the closed form and from
explicit construction, $L = 24$, seven domain sizes from 64 to 3072 solutes.}
\label{tab:gamma}
\small
\begin{tabular}{lrr}
\toprule
& $e_{\rm bulk}$ (eV per solute) & $\gamma$ (eV per bond) \\
\midrule
closed form, Eq.~\eqref{eq:gamma} & $-0.2835$ & $0.1250$ \\
as constructed & $-0.2394$ & $0.1557$ \\
after relaxation & $-0.2802$ & $0.1281$ \\
\bottomrule
\end{tabular}
\end{table}

The relaxed construction agrees with the closed form to 1.2\% in $e_{\rm bulk}$
and 2.5\% in $\gamma$, and the agreement deserves more than a note of numerical
consistency. Eq.~\eqref{eq:gamma} is a continuum statement: it asserts that the
excess energy of a two-phase configuration is strictly proportional to interface
\emph{area}, with a coefficient independent of size or shape. That the same
coefficient reproduces domains of 64 solutes --- objects perhaps a nanometre
across, in which a majority of atoms lie at the surface --- shows that the area
scaling on which the whole continuum treatment rests does not fail at the
nanoscale for this system. The reason is visible in Eq.~\eqref{eq:Ecgamma}: the
cancellation of the bulk terms is exact rather than asymptotic, so no
large-domain limit is being invoked. This is what licenses the use of a
Cahn--Hilliard functional to describe inclusions of the size that natural
sphalerite actually contains. The as-constructed value is higher because the layered
ordering rule cannot be satisfied exactly at a surface: the constructed domains
fall 1.5 to 5.4\% short of the ideal heterovalent bond count, a shortfall that
diminishes with size as the surface fraction falls, and relaxation recovers most
of it. Residuals of the two-term fit are 5.1~eV RMS over an energy range of
450~eV; adding an $n^{1/3}$ edge term reduces this only to 4.92~eV, so the
two-term description is adequate over this range.

\subsection{The cost of splitting a domain}
\label{sec:splitting}

The energy required to divide one domain into two halves is
$\Delta E = \gamma\left[2A(n/2) - A(n)\right]$, evaluated with the relaxed
$\gamma$ and the measured areas:

\begin{center}
\small
\begin{tabular}{rrrrr}
\toprule
$n$ & $A(n)$ & $2A(n/2)$ & $\Delta E$ (eV) & $\Delta E / k_BT$ at 573~K \\
\midrule
128  & 414  & 528  & 14.6 & 296 \\
256  & 642  & 828  & 23.8 & 483 \\
512  & 1040 & 1284 & 31.3 & 633 \\
1024 & 1640 & 2080 & 56.4 & 1141 \\
2048 & 2598 & 3280 & 87.4 & 1769 \\
\bottomrule
\end{tabular}
\end{center}

Every value is positive, as it must be since dividing a domain increases
interfacial area, and every value is enormous relative to thermal energy. This
is the quantitative form of the problem the present paper addresses: bulk
configurational thermodynamics opposes a population of separate domains by
hundreds of $k_BT$, yet such populations are what natural sphalerite frequently
shows. The resolution cannot be thermodynamic.

\subsection{Two sampled estimates that should not be used}
\label{sec:gammafailed}

Two finite-temperature attempts to measure $\gamma$ preceded the closed form and
both failed. They are recorded because the failure modes are diagnostic of a
general difficulty, namely that $\gamma$ is a small difference between large
numbers and sampled configurations carry both energy error and shape noise.

The first relaxed domains isothermally at 573~K, where the acceptance ratio is
of order $10^{-5}$, and left the $n = 1280$ domain at $-145$~eV against an
independently known equilibrium bound of $-298$~eV. The second used replica
exchange and reached $-228$~eV, still 24\% short, and additionally returned a
three-cluster configuration at $n = 1280$, so the measured area belonged to
three domains rather than one. Its diagnostic failure was unmistakable: the
fitted splitting cost at $n = 640$ came out \emph{negative}, which is
impossible.

The closed form avoids both difficulties by not sampling at all. Its verification
in Section~\ref{sec:gammaverify} uses construction and zero-temperature
relaxation, in which the domain outline is imposed rather than inferred.

%==============================================================================
\section{Two competing timescales}
\label{sec:timescales}

Sections~\ref{sec:bridge} and \ref{sec:gamma} establish that a two-phase
configuration carries an interfacial penalty of order hundreds of $k_BT$ per
domain boundary, so a population of separate domains is strongly disfavoured
thermodynamically. \citet{yu2026b} shows the equilibrium state to be a single
condensed domain. Natural sphalerite nonetheless frequently contains many. This
section sets out why, in a form that requires no simulation.

\subsection{The two rates}

Coarsening proceeds by transport: material must leave the smaller domains and
arrive at the larger ones, across a matrix in which every step requires breaking
heterovalent bonds. For diffusion-limited coarsening of domains of radius $R$,
the Lifshitz--Slyozov--Wagner result \citep{lifshitz1961} gives
$R^3 - R_0^3 \propto D t$, so the
time to grow a domain appreciably is
\begin{equation}
\tau_{\rm coarsen}(T) \;\sim\; \frac{R^3\, k_B T}{\gamma\, D(T)\, \Omega},
\qquad
D(T) = D_0 \exp\!\left(-\frac{E_a}{k_B T}\right),
\label{eq:taucoarsen}
\end{equation}
with $\Omega$ the volume per cation site. The interfacial energy appears in the
denominator because it is the driving force: a larger $\gamma$ makes coarsening
faster, not slower.

Cooling supplies the competing rate. What matters is not the absolute
temperature drop but the time over which the mobility changes appreciably, since
that is the interval within which the system either coarsens or does not.
Differentiating the Arrhenius form, the mobility falls by a factor $e$ over a
temperature interval $k_BT^2/E_a$, so
\begin{equation}
\tau_{\rm cool}(T) \;\sim\; \frac{k_B T^2}{E_a\, \lvert \dot{T} \rvert}.
\label{eq:taucool}
\end{equation}

The essential asymmetry is now visible. As temperature falls,
$\tau_{\rm coarsen}$ grows \emph{exponentially}, through the Arrhenius factor,
while $\tau_{\rm cool}$ falls only as $T^2$. Whatever their ordering at high
temperature, the two must cross, and they cross once.

\subsection{The arrest temperature}

Setting $\tau_{\rm coarsen}(\Tf) = \tau_{\rm cool}(\Tf)$ and taking logarithms
gives the arrest condition
\begin{equation}
\frac{E_a}{k_B \Tf}
 = \ln\!\left[ \frac{\gamma\, D_0\, \Omega\, \Tf}
                    {R^3\, E_a\, \lvert \dot{T} \rvert} \right].
\label{eq:arrest}
\end{equation}
Below $\Tf$ the microstructure is frozen: coarsening would still lower the free
energy, and by a large margin, but there is no longer time for it. The
configuration that survives to be observed is the one the system happened to
have reached when transport ceased.

Equation~\eqref{eq:arrest} is transcendental in $\Tf$, but its structure is more
informative than its solution. The left-hand side is large --- with
$E_a \gtrsim 2\lambda \approx 0.6$~eV and $\Tf$ in the ore-forming range,
$E_a/k_B\Tf$ lies between 9 and 17 --- while the right-hand side is a logarithm.
A logarithm cannot vary quickly, so $\Tf$ is strongly buffered against
everything inside it. Differentiating implicitly, with $X = E_a/k_B\Tf$,
\begin{equation}
\frac{\mathrm{d}\Tf}{\Tf} = -\frac{\mathrm{d}\ln \lvert\dot{T}\rvert}{X - 1}.
\label{eq:tfsens}
\end{equation}
A tenfold increase in cooling rate raises $\Tf$ by 14\% at 400~K, 21\% at
573~K and 30\% at 800~K. The arrest temperature is therefore a weak function of
cooling history, and the same is true of its dependence on $D_0$, $\gamma$ and
$R$, all of which enter only logarithmically.

This is a useful robustness rather than a limitation. It means that the
existence of arrest, and its approximate location, do not depend on the
quantities this work cannot pin down: an uncertainty of an order of magnitude in
$D_0$ shifts $\Tf$ by less than a quarter.

\subsection{What is frozen in}

The domain size at arrest follows by evaluating how far coarsening has run when
the two timescales cross. Substituting $\tau_{\rm cool}(\Tf)$ into the
Lifshitz--Slyozov growth law,
\begin{equation}
R_{\rm f}^3 \;\sim\; \frac{\gamma\, D(\Tf)\, \Omega}{k_B \Tf}\,
                     \tau_{\rm cool}(\Tf)
             \;\sim\; \frac{\gamma\, D_0\, \Omega}{E_a\, \lvert\dot{T}\rvert}
                      \, \Tf \exp\!\left(-\frac{E_a}{k_B \Tf}\right).
\label{eq:Rf}
\end{equation}
Because $\Tf$ itself varies only logarithmically with $\lvert\dot{T}\rvert$, the
explicit factor dominates and
\begin{equation}
\boxed{\;R_{\rm f} \;\propto\; \lvert\dot{T}\rvert^{-n}, \qquad n \to 1/3 \;}
\label{eq:scaling}
\end{equation}
in the limit where the drift in $\Tf$ is neglected, with $n$ reduced below
$1/3$ once it is not, because faster cooling arrests at a higher temperature
where the mobility is larger and thus partly compensates. The phase-field
simulations of Section~\ref{sec:pf} return $n = 0.259$, consistent with a
one-third law softened by exactly that drift.

Equation~\eqref{eq:scaling} is the central predictive statement of this paper.
It says that the size of the inclusions a sphalerite grain carries is set by how
fast it cooled, through a power law whose exponent is a property of
diffusion-limited coarsening rather than of any parameter this work has had to
assume. Section~\ref{sec:geospeedometer} takes up its inversion.

\section{Phase-field simulation}
\label{sec:pf}

The analysis of Section~\ref{sec:timescales} predicts arrest and a power-law
dependence of the frozen domain size on cooling rate, but says nothing about the
morphology that results. This section supplies it, subject throughout to the
division of labour set out in Section~\ref{sec:kappacaveat}: the phase-field
calculation gives the pathway and the topology, not the absolute scale.

\subsection{Implementation}

Equation~\eqref{eq:ch} is integrated on a periodic square grid by a
semi-implicit spectral scheme. The fourth-order linear term is diagonal in
Fourier space and is treated implicitly, the nonlinear bulk term explicitly:
\begin{equation}
\hat{c}^{\,n+1} = \frac{\hat{c}^{\,n} - \Delta t\, M k^2\,
                        \mathcal{F}\!\left[f'(c^n)\right]}
                       {1 + \Delta t\, M \kappa k^4}.
\label{eq:scheme}
\end{equation}
An explicit scheme would require $\Delta t < C\,\Delta x^4/(M\kappa)$, which at
the resolutions needed here is unusable.

Two implementation choices require comment because both were arrived at by
first getting them wrong.

\paragraph{The gradient coefficient is set from the resolvable wavelength.}
Linear stability of Eq.~\eqref{eq:ch} gives a fastest-growing wavelength
$\lambda = 2\pi\sqrt{2\kappa/|f''|}$. At the physical $\kappa$ this is of order
the lattice parameter and no grid resolves it: with $\kappa$ from
Eq.~\eqref{eq:gammakappa} the fastest-growing mode spans one cell, the interface
is unresolved, and the calculation describes numerical diffusion rather than
decomposition. We therefore set $\kappa = |f''|\lambda_{\rm target}^2/8\pi^2$
with $\lambda_{\rm target} = 12$ cells. This is the widening announced in
Section~\ref{sec:kappacaveat}, here made quantitative.

\paragraph{The timestep is set from the target amplification, not from a CFL
condition.} Linearising Eq.~\eqref{eq:scheme}, a mode $k$ is amplified per step
by $A(k) = (1 - \Delta t M k^2 f'')/(1 + \Delta t M \kappa k^4)$, so growth
requires $-f'' > \kappa k^2$, which is the stability condition and does not
involve $\Delta t$. What $\Delta t$ controls is the amount of growth per step.
Setting it from an explicit-scheme CFL condition gives $A - 1 \sim 3\times
10^{-4}$, so the unstable mode needs some 7000 steps to grow tenfold while the
run has only a few thousand at that mobility; the result looks like diffusive
smoothing, because that is what it is. We instead solve for the $\Delta t$ that
delivers a chosen amplification, $A = 1.02$, of the fastest-growing mode. Taking
a large $\Delta t$ is precisely what the implicit treatment of the stiff term
exists to permit.

\subsection{The morphological sequence}

Figure~\ref{fig:sequence} follows a cooling run from 1800 to 500~K. The
composition starts uniform at $c_0 = 0.30$ with fluctuations of amplitude
$2\times10^{-3}$, which places it inside the unstable region of
Section~\ref{sec:spinodal}.

\begin{figure}[htbp]
\centering
\includegraphics[width=\textwidth]{fig_sequence_32000.pdf}
\caption{Morphological evolution of a slowly cooled grain. Left to right:
initial fluctuations; spinodal decomposition, which raises the composition
contrast from $2\times10^{-3}$ to 0.308 and establishes two phases; and
progressive coarsening, in which the domain count falls and the surviving
domains grow. Panels are individually scaled, since a common scale set by the
seeding noise would compress every later panel. Lengths are in grid units, not
physical units: see Section~\ref{sec:kappacaveat}.}
\label{fig:sequence}
\end{figure}

The sequence has four stages, and they are the stages the introduction promised.
The field is initially homogeneous to within the seeding noise. Decomposition
then occurs spinodally: the composition contrast rises by more than two orders
of magnitude within the first temperature interval, and a two-phase structure
appears with a wavelength close to the linear-stability prediction. Coarsening
follows, the domain count falling from many to few as material transfers from
smaller domains to larger. Finally the microstructure ceases to evolve: the
last panels differ from one another far less than the early ones, not because
coarsening has completed --- a single domain would be the completed state ---
but because the mobility has fallen too far for it to continue.

\subsection{Cooling rate and the frozen microstructure}

Figure~\ref{fig:rates} compares three cooling schedules. Since $\Delta t$ tracks
the mobility, specifying the schedule as a number of steps per kelvin fixes the
diffusive progress per kelvin, which is the dimensionless ratio
Section~\ref{sec:timescales} identifies as controlling. A larger step count is
slower cooling.

\begin{figure}[htbp]
\centering
\includegraphics[width=\textwidth]{fig_cooling_rates.pdf}
\caption{Frozen microstructures and coarsening histories at three cooling
rates. Upper row: the final state, with the domain count. Lower panel: domain
scale against temperature, decreasing to the right. Every curve flattens as the
mobility collapses --- that flattening is arrest --- and the slower the cooling,
the further coarsening runs before it happens.}
\label{fig:rates}
\end{figure}

\begin{table}[htbp]
\centering
\caption{Frozen microstructure against cooling rate, $256\times256$ grid,
$\lambda_{\rm target} = 12$ cells. Domain scale is from the first moment of the
structure factor, in grid units.}
\label{tab:rates}
\small
\begin{tabular}{rlrr}
\toprule
cooling steps & rate & domain scale $L$ & domains \\
\midrule
 2{,}000 & fast   & 18.2 & 160 \\
 8{,}000 & medium & 26.5 &  76 \\
32{,}000 & slow   & 37.3 &  33 \\
\bottomrule
\end{tabular}
\end{table}

The frozen domain scale follows $L \propto (\text{steps})^{0.259}$, that is
$L \propto |\dot{T}|^{-0.259}$, against the $1/3$ of Eq.~\eqref{eq:scaling}
softened by the drift in $\Tf$ that Eq.~\eqref{eq:tfsens} describes. Analysis
and simulation therefore agree on the exponent to within the accuracy either
supports.

\subsection{A caution on what agreement demonstrates}

The exponent proved insensitive to errors it should have been sensitive to, and
the episode is worth recording because it bears on how such agreement should be
read. Two earlier versions of this calculation returned $0.242$ and $0.256$ for
the same quantity. In the first the interface was unresolved; in the second the
free energy was internally inconsistent, the analytic derivative $f'(c)$
disagreeing with $f''(c)$ by a factor of 150 at the working composition because
a term had been dropped from the differentiation. In neither case did
decomposition actually occur --- the composition contrast decayed rather than
grew --- yet both produced plausible pictures of a coarsening pattern and an
exponent within 7\% of the correct one.

The reason is that the exponent reflects the transport-limited kinetics of
whatever structure is present, and pattern coarsening under pure diffusion obeys
a similar law to domain coarsening under phase separation. Agreement in the
exponent was therefore not evidence that the mechanism was right. What exposed
the error was tracking the composition contrast, which must grow by orders of
magnitude if decomposition is occurring and did not. We now assert internal
consistency of $f$, $f'$ and $f''$ against finite differences whenever the free
energy is constructed, and report the contrast alongside the domain scale in
every figure.

\section{Direct coarsening in the lattice model}
\label{sec:latticecoarsen}

The phase-field calculation cannot supply an absolute length, for the reason set
out in Section~\ref{sec:kappacaveat}. This section supplies one, by coarsening
the lattice model directly. Here the interface is what it physically is, one
bond wide, and nothing is inflated; what the calculation gives up in reach it
gains in meaning, since sizes come out in atoms and hence in nanometres.

\subsection{Protocol, and two design errors that had to be corrected}

A dispersed solute distribution is held at fixed temperature and the domain
population followed. Both corrections below were forced by the data and are
recorded because either would silently corrupt the exponent.

\paragraph{The measurement must exclude monomers.} At 1800~K some 90\% of
clusters are single solutes evaporated from the domains. Including them in a
mean radius makes that quantity a measure of the free-solute concentration
rather than of the domains: it \emph{falls} while the domains grow, because a
growing domain sheds more monomers. A first version did include them and
returned exponents of 0.007, 0.024 and 0.165 at 1800, 1600 and 1400~K --- not a
coarsening exponent but an artefact of the mixing. We now average over clusters
of at least four sites and report the monomer fraction separately, so that
evaporation is visible rather than hidden.

\paragraph{The measurement must end before the solute budget does.} At
$x = 1$~at.\% in an $L = 30$ box the 1080 solutes condensed into a single domain
within 3000 of 60{,}000 sweeps, so more than nine tenths of the run measured a
structure that had stopped evolving, dragging the fitted exponent towards zero.
That the equilibrium state is a single domain is the central result of
\citet{yu2026b}; it follows that a coarsening measurement must be made before
exhaustion, and therefore that enough solute must be supplied for exhaustion to
take a long time. At $x = 4$~at.\% the coarsening window extends over the whole
run.

The composition is thus far above natural indium tenors. It is chosen to make
the exponent measurable, and the exponent is a property of the transport-limited
kinetics rather than of the concentration; the absolute sizes reported are those
of the simulated domains, not a prediction of natural inclusion sizes.

\subsection{Results}

In an $L = 24$ box (55{,}296 cation sites, a cube 13.0~nm on a side) at 1600~K,
the domain count falls from 173 to 4 over 20{,}000 sweeps while the mean domain
radius grows from 0.44 to 1.66~nm. These are absolute figures: the simulated
domains are nanometric, in a box whose edge is 13~nm, with no adjustable length
scale anywhere in the calculation.

The fitted exponent is $n = 0.150 \pm 0.002$, and it is stable across the run:
taking successive decades of sweeps separately gives 0.176, 0.137, 0.156 and
0.122. It is not an artefact of the fitting window.

\subsection{Why the exponent is below one third}

The Lifshitz--Slyozov--Wagner value is $1/3$ and the measurement is half that,
so the difference requires an explanation rather than a footnote.

The obvious suspicion is that solute continues to arrive from the dispersed
population, so that new small domains are counted while the large ones grow,
depressing the mean. That can be tested: if the counted domains hold a
conserved amount of solute, then $N R^3$ is constant and the exponent obtained
from the domain count must agree with the one from the radius. It does. The
product $N R^3$ varies between 13 and 19 with no systematic drift, the count
gives $N \sim t^{-0.400}$ and hence $R \sim t^{+0.133}$, and the radius gives
$R \sim t^{+0.148}$. Mass is conserved among the domains and the two routes
agree, so continuing nucleation is not the explanation.

Three possibilities remain, and we cannot presently distinguish them.

The domains may be too small for the asymptotic law. Lifshitz--Slyozov assumes
domains large compared with the interface width and dilute enough that they do
not interact; here the domains are a few nanometres across with interfaces one
bond wide, and at $x = 4$~at.\% they are not dilute. Pre-asymptotic exponents
below $1/3$ are commonly reported in such regimes.

The transport may not be purely diffusive. The Monte Carlo dynamics conserves
composition but has no vacancy mechanism: material moves by direct exchange,
which is not the mechanism operating in a real crystal and need not give the
same exponent.

The run may not have reached the asymptotic regime. Twenty thousand sweeps
covers less than three decades in time, and the domain count falls to four, so
the late window is statistically thin.

What the measurement establishes firmly is the \emph{scale}: coarsening in this
model produces domains of one to two nanometres over the accessible window, in
a box of 13~nm, with the interface unmodified. That is the absolute figure the
phase-field calculation could not provide, and it is the figure that
Section~\ref{sec:geospeedometer} compares with observation. The exponent from
this route should be treated as a lower bound on the asymptotic value rather
than as a determination of it.

\section{Cross-check between the two methods}
\label{sec:crosscheck}

Sections~\ref{sec:pf} and \ref{sec:latticecoarsen} describe the same physical
process with representations that could hardly be more different. The
phase-field calculation carries a continuous order parameter on a grid whose
interface has been deliberately widened to span a dozen cells; the lattice
calculation carries discrete species on the cation sublattice with an interface
one bond wide. Their only shared input is the free energy, and even that enters
differently: as a smooth double well in one, as an explicit sum over
configurations in the other.

The coarsening exponent is the quantity both can compute, and it is therefore
the natural cross-check. If the two agree, the exponent is a property of the
transport-limited kinetics rather than of either representation --- which is the
premise on which Section~\ref{sec:geospeedometer} rests, since the inversion
uses the exponent and nothing else.

We compare three determinations:
\begin{itemize}\itemsep2pt
\item the Lifshitz--Slyozov--Wagner result for diffusion-limited coarsening,
      $n = 1/3$, softened by the drift in $\Tf$ of Eq.~\eqref{eq:tfsens} to
      somewhat below that value;
\item the phase-field cooling runs, which give $n = 0.259$;
\item direct lattice coarsening, which gives $n = 0.150 \pm 0.002$
      (Section~\ref{sec:latticecoarsen}).
\end{itemize}

The three do not agree. Theory gives $1/3$, the phase-field runs $0.259$, and
the lattice $0.150$; the natural end members of Section~\ref{sec:geospeedometer}
imply something near $0.45$. The spread is a factor of three, and it would be
dishonest to present the two simulation values as mutual corroboration when they
differ by 70\%.

What the comparison does establish is a hierarchy of confidence. Both
simulations reproduce the qualitative law --- domains grow as a power of time,
and the frozen size falls with cooling rate --- and both do so with
representations that share almost nothing. The exponent itself is evidently
sensitive to things the two treat differently: the phase-field interface is
twelve cells wide and the lattice interface one bond, the phase-field dynamics
is continuum diffusion and the lattice dynamics direct exchange, and the lattice
domains are small enough that the asymptotic law need not hold. Each of those
would be expected to move the exponent, and Section~\ref{sec:latticecoarsen}
argues that the lattice value should be read as a lower bound.

The inversion of Section~\ref{sec:geospeedometer} is therefore presented as an
ordering rather than a calibration. Ranking deposits by cooling rate requires
only that the exponent be positive and roughly constant, which all four
determinations support. Returning a rate in kelvin per year would require
knowing it to better than a factor of two, which none of them presently does.

Two cautions bear on how such agreement should be weighed, and both were learned
the hard way.

The first is the episode of Section~\ref{sec:pf}: two versions of the
phase-field calculation in which decomposition did not occur at all returned
exponents of $0.242$ and $0.256$, within 7\% of the correct value. Coarsening of
a diffusive pattern obeys a similar law to coarsening of a phase-separated
microstructure, so the exponent alone does not identify the mechanism. We
therefore report the composition contrast alongside every exponent: in the
lattice model the analogous check is the fraction of solute in monomers, which
must fall towards zero if domains are genuinely forming.

The second concerns temperature. Coarsening at 573~K is unobservable in the
lattice model, since the acceptance ratio there is of order $10^{-5}$. The
lattice runs are therefore conducted at 1400 to 1800~K, on the assumption that
the exponent is a property of the kinetics rather than of the temperature. That
assumption is testable within the calculation, by measuring $n$ at several
temperatures and asking whether it moves, and we report the spread rather than a
single value.

%==============================================================================
\section{Inversion: nanoinclusions as a geospeedometer}
\label{sec:geospeedometer}

Equation~\eqref{eq:scaling} states that the frozen inclusion size depends on the
cooling rate through a power law. Read forwards it predicts a microstructure
from a thermal history. Read backwards it does the opposite, and that is the
more useful direction, because inclusion sizes are measurable in a thin section
whereas cooling rates are not.

\subsection{The relation to be inverted}

Inverting Eq.~\eqref{eq:scaling},
\begin{equation}
\lvert \dot{T} \rvert \;=\; \lvert \dot{T} \rvert_{\rm ref}
      \left( \frac{R_{\rm f}}{R_{\rm ref}} \right)^{-1/n},
\label{eq:invert}
\end{equation}
so a single calibration point $(R_{\rm ref}, |\dot{T}|_{\rm ref})$ plus the
exponent converts an observed size into a rate. Because $n \approx 1/3$, the
inversion carries a factor of three amplification: a size measured to within a
factor of two yields a rate to within a factor of eight. That is coarse, but
cooling rates in ore systems span ten orders of magnitude, so a factor of eight
still discriminates strongly between the deposit classes considered below.

\subsection{Natural end members}

Indium-bearing sphalerite occurs in deposit types whose cooling rates differ by
some eight orders of magnitude, and the inclusion textures reported from them
differ correspondingly. Table~\ref{tab:endmembers} sets the two against each
other.

\begin{table}[htbp]
\centering
\caption{Cooling-rate end members for In-bearing sphalerite and the inclusion
textures reported from them. Rates are order-of-magnitude estimates for the
deposit class, not measurements on the specific samples in which the textures
were observed; the comparison is between populations, not between paired
observations.}
\label{tab:endmembers}
\small
\begin{tabular}{llll}
\toprule
deposit class & cooling rate (K\,yr$^{-1}$) & inclusion texture & scale \\
\midrule
VMS, seafloor massive sulphide & $10^{0}$--$10^{2}$ &
  lattice distortion, ultrafine domains & 5--50\,nm \\
magmatic-hydrothermal, skarn & $10^{-3}$--$10^{-1}$ &
  sub-micron exsolution & 0.1--1\,\textmu m \\
MVT, basinal brine & $10^{-6}$--$10^{-4}$ &
  optically visible roquesite & 1--50\,\textmu m \\
\bottomrule
\end{tabular}
\end{table}

The qualitative ordering is exactly as Eq.~\eqref{eq:scaling} requires: the
faster the system cooled, the finer the inclusions it carries. The
quenched seafloor systems, in which hydrothermal fluid meets cold seawater, are
where indium most often presents as an apparent solid solution in LA-ICP-MS
depth profiles --- smooth ablation signals with no discrete spikes --- and where
transmission electron microscopy nonetheless resolves nanoscale heterogeneity
\citep{cook2009,frenzel2016}. At the other extreme, roquesite is reported as
discrete micrometric grains exsolved from sphalerite \citep{moura2007}.

\subsection{A quantitative discrepancy, and what it may mean}

The ordering agrees. The magnitude does not, and we report the disagreement
rather than the agreement, because it is the more informative of the two.

Across the full range of Table~\ref{tab:endmembers}, cooling rate spans eight
decades and inclusion size spans four, which implies
\begin{equation}
n_{\rm obs} \;=\; \frac{\Delta \log_{10} R}{\Delta \log_{10} |\dot{T}|}
           \;\approx\; 0.50 .
\label{eq:nobs}
\end{equation}
Taking the classes pairwise rather than end to end gives $0.43$ and $0.45$,
consistent with each other and with Eq.~\eqref{eq:nobs}. The theoretical and
simulated values are $1/3$ and $0.259$. The natural data therefore imply an
exponent larger by roughly a factor of $1.7$, and the two pairwise estimates
agree closely enough that this is unlikely to be scatter.

Several readings are possible and we cannot presently choose between them.

\emph{The calibration may be miscast.} Table~\ref{tab:endmembers} pairs
deposit-class cooling rates with textures reported from different samples, not
from the same ones. If the fastest-cooled samples in which nanodomains are
resolved are not the fastest-cooling members of their class, or if optically
visible roquesite is preferentially reported from the slowest members of theirs,
the apparent range of both axes is inflated in a correlated way. Paired
measurements on single samples would settle this and do not appear to exist.

\emph{The coarsening mechanism may differ at the extremes.} The
Lifshitz--Slyozov exponent of $1/3$ assumes coarsening limited by bulk
diffusion. If transport along dislocations or grain boundaries contributes in
the slowly cooled systems, the effective exponent rises; interface-limited
coarsening gives $1/2$ exactly, which is close to what the natural data imply.
The discrepancy may therefore be a signature that the slow end members coarsen
by a different mechanism from the one modelled here.

\emph{The arrest temperature may drift more than assumed.}
Equation~\eqref{eq:tfsens} treats $E_a$ as constant. If the activation energy
itself falls as domains grow --- plausible, since a solute leaving a large domain
breaks fewer bonds per unit of interface removed --- then $\Tf$ drifts further
than estimated and the effective exponent changes.

We regard the third as least likely, because it would move the exponent in the
wrong direction, and the first two as both plausible. What the comparison
establishes is that the model predicts the correct ordering across eight decades
of cooling rate, and predicts a magnitude that is too small by a factor of about
$1.7$; and that the resolution lies either in better-paired observations or in
coarsening mechanisms this work does not include.

\subsection{What the tool can and cannot do}

Stated plainly, so that it is not overread.

It \emph{can} rank deposits by cooling rate from inclusion textures, which is
what the agreement in ordering supports. A sphalerite carrying optically visible
roquesite cooled far more slowly than one whose indium is invisible to
LA-ICP-MS, and the difference is many orders of magnitude rather than a factor
of a few.

It \emph{cannot} yet return a cooling rate in kelvin per year from a measured
inclusion size, for two reasons. The absolute calibration requires a diffusion
coefficient this work does not have (Section~\ref{sec:mobility}), and the
exponent disagrees with the natural data by a factor whose origin is unresolved.
Until both are settled, the inversion should be read as a relative one.

\subsection{The rate-limiting step}

One mineralogical observation supports the framework and deserves recording,
because it bears on the activation energy that controls arrest. Copper diffuses
anomalously rapidly in sphalerite, plausibly by an interstitial mechanism; this
is the basis of the classic account of chalcopyrite disease
\citep{barton1987}. Indium is a different case: In$^{3+}$ is large, highly
charged, and cannot move without either a vacancy or a compensating counter-ion
moving with it, so its diffusivity must be far below that of copper and below
zinc self-diffusion as well.

The consequence for this paper is direct. Coarsening of a Cu--In domain requires
both species to move, so it proceeds at the rate of the slower, and the arrest
condition of Eq.~\eqref{eq:arrest} should be evaluated with the activation
energy for indium rather than an average. Arrest therefore occurs at a higher
temperature than a copper-based estimate would suggest, and the microstructure
freezes earlier. The nanoscale inclusions that natural sphalerite carries are,
on this reading, a record of indium having stopped moving while the rock was
still hot.

\section{Discussion}
\label{sec:discussion}

\subsection{A change in what a texture means}

Trace-element textures in sphalerite are conventionally read as a record of
\emph{composition}: what the fluid supplied, and what the crystal could
accommodate. The argument of this paper makes them a record of something else
as well. If inclusions are frozen coarsening products, their size is set not by
how much indium was available but by how long the system had to reorganise it
before transport stopped. Two grains with identical indium contents, grown from
identical fluids, will carry different textures if they cooled differently.

That reframing has a corollary worth stating explicitly, because it inverts an
assumption. A smooth LA-ICP-MS depth profile has often been taken as evidence
for solid solution. On the present account it is evidence for \emph{fast
cooling}: the indium is aggregated, as \citet{yu2026b} requires, but into
domains too fine for the ablation volume to resolve. The absence of spikes
constrains the thermal history, not the state.

\subsection{Ranking deposits by thermal history}

Section~\ref{sec:geospeedometer} establishes that the model reproduces the
ordering of inclusion size across deposit classes spanning eight orders of
magnitude in cooling rate, while disagreeing on the magnitude by a factor of
about $1.7$. The ordering is the part that is usable now.

Its practical form is comparative rather than absolute. Given two sphalerite
samples, the one with coarser indium-bearing inclusions cooled more slowly, and
by a margin that the power law converts into orders of magnitude rather than
factors. Within a single deposit, systematic variation of inclusion size across
a district would map relative cooling rates, and therefore proximity to the heat
source, in a way that bulk indium tenor cannot: tenor reflects supply, texture
reflects time.

We are careful not to overstate what follows. Turning this into a
geothermometer proper --- returning a cooling rate in kelvin per year --- would
require an absolute diffusion coefficient, which Section~\ref{sec:mobility}
explains this work does not have, and an exponent known to better than a factor
of two, which Section~\ref{sec:crosscheck} shows none of the four
determinations presently provides. Both are obtainable; neither is obtained
here.

\subsection{Consequences for metallurgical recovery}

Indium is recovered almost entirely as a by-product of zinc refining, so its
recovery efficiency depends on how it is bound in the sphalerite fed to the
smelter. The state of that binding is precisely what this work argues is set by
thermal history.

The two limiting cases behave differently. Indium dispersed as nanometric
domains within the sphalerite lattice, the expected product of rapid cooling,
will follow zinc through dissolution and report to the same process streams;
its recovery is a question of hydrometallurgical circuit design. Indium
concentrated in discrete micrometric roquesite, the product of slow cooling,
is a separate mineral phase with its own dissolution behaviour, and may be
recoverable --- or lost --- by physical separation before leaching.

The practical suggestion is therefore that deposit classes be expected to
differ systematically in indium deportment, and along an axis that is
predictable in advance from cooling history rather than discoverable only by
metallurgical testing. Quenched seafloor systems should present indium in the
lattice-bound form; slowly cooled basinal systems should present a greater
fraction as discrete inclusions. We advance this as a prediction the framework
makes, not as a result it establishes, since we have tested it against
published textures rather than against recovery data.

\subsection{The generality of the argument}

Little in Sections~\ref{sec:timescales} and \ref{sec:geospeedometer} is specific
to indium. The ingredients are a solute whose equilibrium state is condensed, a
transport process that is thermally activated, and a host that cools. Any
trace element meeting those conditions should freeze in a coarsening
microstructure whose scale records the cooling rate, with the same power law and
a different prefactor.

What \emph{is} specific to indium is the reason its equilibrium state is
condensed, which is the charge-compensation physics of \citet{yu2026}: the
coupled substitution of Cu$^+$ and In$^{3+}$ for two Zn$^{2+}$ carries a pair
binding energy fixed by the dielectric constant and an order of magnitude above
$k_BT$. Elements substituting isovalently have no such driving force. The
framework should therefore transfer to other heterovalent couples --- Ag--Sb,
Cu--Ga --- and not to isovalent substituents such as Fe or Cd, whose
distribution is governed by different physics.

\subsection{On the disagreement between methods}

Section~\ref{sec:crosscheck} reports that four determinations of the coarsening
exponent span a factor of three, and it would be convenient to attribute the
spread to numerical detail. We think the more likely reading is that the
exponent is genuinely different in the different regimes, and that this is
informative rather than merely inconvenient.

The lattice value of $0.150$ comes from domains a few nanometres across with
atomically sharp interfaces, well outside the regime in which the
Lifshitz--Slyozov derivation holds. The natural end members imply something near
$0.45$, which is close to the $1/2$ of interface-limited rather than
diffusion-limited coarsening --- and interface limitation is plausible for
domains whose interface is one bond wide. The phase-field value of $0.259$ falls
between, from a calculation whose interface is deliberately widened towards the
continuum limit.

Read that way, the spread would be a crossover: from interface-limited
coarsening at the atomic scale towards diffusion-limited coarsening as domains
grow. That reading is testable, by measuring the exponent as a function of
domain size within a single simulation rather than comparing across
simulations that differ in several respects at once.

\subsubsection*{The crossover hypothesis, tested and not supported}

We have since carried out that test, and report it here because it removes one
of the candidate explanations rather than because it succeeded.

If $n$ depends on domain size relative to interface width, then runs at
different interface widths should give curves that differ when plotted against
$L$ but coincide when plotted against $L/w$. Four isothermal Cahn--Hilliard
runs were made on a $512^2$ grid at $w = 1.0$, 1.5, 2.25 and 3.4 cells, with the
step count allocated as $w^{-1/n}$ so that each reached the largest domain size
its width allowed: the narrow-interface run reached $L/w = 70$, the widest
stopped at 18. The local exponent was extracted with a sliding window whose
width had been calibrated beforehand against synthetic traces carrying a known
crossover, so that its bias and scatter were established independently of the
data under test.

The curves do not collapse. Rescaling the abscissa by $w$ reduces the scatter
between them only from 0.011 to 0.008, a ratio of 0.75, where collapse would
require a ratio well below one. A more direct statement of the same result: at
a common $L/w$ of 4.4 the four runs give $n = 0.203$, 0.205, 0.212 and 0.208,
agreeing to 0.004 --- but at that point their absolute $L$ differs by a factor
of 3.4, so $n$ tracks neither variable.

Two further features contradict the hypothesis outright. The exponent
\emph{rises} through every run, from $0.207 \pm 0.004$ to $0.267 \pm 0.006$,
where the hypothesis requires a fall from about $1/2$ towards $1/3$. And every
value lies far below both candidates: no run approaches either 0.5 or 0.333 at
any stage.

What the exponent does track is how far the run has progressed. The value
reached at the end of each run correlates with the logarithm of its step count
at $r = 0.99$, which is what would be expected if all four are approaching a
common asymptote from below and have simply gone different distances towards
it. On that reading the phase-field exponent of 0.259 is not a regime-specific
value at all but a transient one, and the asymptote --- which may well be
$1/3$ --- lies beyond the reach of these runs.

The crossover hypothesis is therefore not supported, and the discrepancy of
Section~\ref{sec:geospeedometer} returns to the explanation offered first
there: that the end-member comparison pairs cooling rates with textures
reported from different samples, and that paired measurements on single
specimens are what would settle it.

\section{Limitations}
\label{sec:limitations}

We set these out in the order in which they bear on the conclusions, and
distinguish throughout between limitations that affect the numbers and
limitations that would affect the argument.

\paragraph{The phase-field interface is artificially widened.} This is stated
in Section~\ref{sec:kappacaveat} before any phase-field result is presented, and
it is absolute: the domain sizes in Figures~\ref{fig:sequence} and
\ref{fig:rates} are set by the grid spacing, not predicted. Inflating $\kappa$
rescales lengths and times together, so the dimensionless conclusions --- the
order of morphologies, the existence of arrest, the sign and approximate
magnitude of the size--rate exponent --- survive; the absolute scale does not,
and comes instead from Section~\ref{sec:latticecoarsen}.

\paragraph{The lattice dynamics has no vacancy mechanism.} Solute moves by
direct exchange, which conserves composition and satisfies detailed balance but
is not how atoms move in a real crystal. Indium in particular is expected to
migrate only with a vacancy or a compensating counter-ion, as
Section~\ref{sec:geospeedometer} discusses, and a mechanism with a different
correlation factor need not give the same coarsening exponent. This is a
plausible contributor to the $0.150$ measured in Section~\ref{sec:latticecoarsen}
and is the reason that value is offered as a lower bound rather than a
determination.

\paragraph{The simulated compositions are far above natural tenors.} The lattice
coarsening runs use 4~at.\% total solute, some two to three orders of magnitude
above the indium content of economically relevant sphalerite, because at natural
dilution the solute budget is exhausted into a single domain before any
coarsening can be measured --- which is itself the result of \citet{yu2026b}.
The exponent is a property of the transport-limited kinetics rather than of the
concentration, but the domains at 4~at.\% are not dilute, and interactions
between them are a further reason the asymptotic law need not hold.

\paragraph{Monte Carlo sweeps are not physical time, and the mobility is not
absolute.} Section~\ref{sec:mobility} sets out the consequence: results are
relative calibrations. A ratio of domain sizes maps onto a ratio of cooling
rates; neither maps onto an absolute value without a measured diffusion
coefficient for Cu or In in sphalerite.

\paragraph{The coarsening exponent is not settled.} Four determinations span a
factor of three (Section~\ref{sec:crosscheck}). Section~\ref{sec:discussion}
suggests they may describe different regimes rather than disagreeing about one,
but that suggestion is untested. Until it is, the inversion supports ranking
and not calibration.

\paragraph{The end-member comparison is not paired.} Table~\ref{tab:endmembers}
sets deposit-class cooling rates against inclusion textures reported from
different samples. Paired measurements --- inclusion size distributions and
independent thermal-history constraints on the same specimens --- would test the
power law properly and do not appear to exist. This is the single most useful
observation the framework calls for.

\paragraph{Iron is absent, and the chemical term is assumed.} Both are inherited
from \citet{yu2026} and discussed there. Fe$^{2+}$ is isovalent with
Zn$^{2+}$ and so does not enter the charge-compensation term, but modifies the
dielectric response that fixes $\lambda$; the chemical term $J_{\alpha\beta}$
remains unconstrained by first-principles calculation, and enters $\gamma$
explicitly through Eq.~\eqref{eq:gamma}.

\paragraph{What none of these affects.} The existence of kinetic arrest does not
depend on any of the above. It follows from two facts that are not in question:
that solute transport is thermally activated, so the coarsening time grows
exponentially as temperature falls; and that cooling proceeds at a rate which
does not. Two such functions cross, and they cross once. Every quantity that
might be wrong --- the diffusion prefactor, the activation energy, the
interfacial energy, the domain size --- enters the arrest condition of
Eq.~\eqref{eq:arrest} inside a logarithm, and Eq.~\eqref{eq:tfsens} shows that
an order of magnitude error in any of them shifts the arrest temperature by less
than a quarter.

What is uncertain is therefore the calibration, not the mechanism. The
inclusions in natural sphalerite are frozen coarsening products; how precisely
their size dates the freezing is what remains to be settled.

%==============================================================================
\section*{Data and code availability}
\label{sec:availability}
Source code, numerical results and figures for this study are openly available
at \url{https://github.com/Ruqing1963/sphalerite-kinetic-arrest} and
archived at \url{https://doi.org/10.5281/zenodo.21995641}. Code is
released under the MIT licence and data under CC~BY~4.0; the manuscript remains
the copyright of the authors. The repository includes the Cahn--Hilliard solver
with Arrhenius mobility used in Section~\ref{sec:pf}, the direct lattice
coarsening of Section~\ref{sec:latticecoarsen}, the thermodynamic integration
that measures $f(c)$ in Section~\ref{sec:fofc}, and the scripts that regenerate
each figure.

Two components are released although they were not run for this work. The
structure generation and least-squares harness for the first-principles
determination of $J_{\alpha\beta}$ is complete and validated against synthetic
energies, so it can be run by anyone with sufficient memory and the parameters
dropped into Eq.~\eqref{eq:gamma} without repeating any simulation reported
here. And the verification suite of \citet{yu2026} is included unchanged, since
every result in this paper depends on the lattice module passing it.

The two preceding papers of the series are archived separately, with their own
code: \citet{yu2026} at \url{https://doi.org/10.5281/zenodo.21880502} and
\citet{yu2026b} at \url{https://doi.org/10.5281/zenodo.21911898}.

\section*{Funding}
This study was funded by the Hezhou Municipal Scientific Research and
Development Program (Project Nos.\ 2024143, 2024141 and 2024104).

%==============================================================================
\begin{thebibliography}{99}
\itemsep2pt

\bibitem[Barton and Bethke(1987)]{barton1987}
Barton, P.B., Bethke, P.M., 1987. Chalcopyrite disease in sphalerite: pathology
and epidemiology. \textit{American Mineralogist} 72, 451--467.

\bibitem[Cahn and Hilliard(1958)]{cahn1958}
Cahn, J.W., Hilliard, J.E., 1958. Free energy of a nonuniform system. I.
Interfacial free energy. \textit{Journal of Chemical Physics} 28, 258--267.

\bibitem[Cook et al.(2009)]{cook2009}
Cook, N.J., Ciobanu, C.L., Pring, A., Skinner, W., Shimizu, M., Danyushevsky,
L., Saini-Eidukat, B., Melcher, F., 2009. Trace and minor elements in
sphalerite: a LA-ICPMS study. \textit{Geochimica et Cosmochimica Acta} 73,
4761--4791.

\bibitem[Frenzel et al.(2016)]{frenzel2016}
Frenzel, M., Hirsch, T., Gutzmer, J., 2016. Gallium, germanium, indium and
other trace and minor elements in sphalerite as a function of deposit type ---
a meta-analysis. \textit{Ore Geology Reviews} 76, 52--78.

\bibitem[Lifshitz and Slyozov(1961)]{lifshitz1961}
Lifshitz, I.M., Slyozov, V.V., 1961. The kinetics of precipitation from
supersaturated solid solutions. \textit{Journal of Physics and Chemistry of
Solids} 19, 35--50.

\bibitem[Moura et al.(2007)]{moura2007}
Moura, M.A., Botelho, N.F., de Mendon\c{c}a, F.C., 2007. The indium-rich
sulfides and rare arsenates of the Sn-In-mineralized Mangabeira A-type granite,
central Brazil. \textit{The Canadian Mineralogist} 45, 485--496.

\bibitem[Yu et al.(2026a)]{yu2026}
Yu, H., Chen, R., Lu, H., 2026a. A statistical thermodynamic lattice model for
Cu--In coupled substitution in sphalerite, version 2.0. \textit{Zenodo}.
\url{https://doi.org/10.5281/zenodo.21880502}

\bibitem[Yu et al.(2026b)]{yu2026b}
Yu, H., Chen, R., Lu, H., 2026b. The equilibrium fate of indium in sphalerite:
large-scale replica-exchange simulation and the absence of a dilute solid
solution. \textit{Zenodo}. \url{https://doi.org/10.5281/zenodo.21911898}

\end{thebibliography}

\end{document}
